\chapter{Conclusiones y perspectivas}

El presente trabajo de tesis representa un esfuerzo integral en el diseño y control de la aeronave EHEKATL BWB. A lo largo de esta investigación, se llevaron a cabo una serie de etapas fundamentales que abordaron desde el diseño aerodinámico hasta la implementación de controladores avanzados, culminando en simulaciones de vuelo en el entorno de X-PLANE. 

Durante la fase de diseño, se utilizó una combinación de herramientas especializadas, como XFLR5, Airoil Maker y Plane Maker, para concebir y evaluar el perfil aerodinámico de la aeronave EHEKATL BWB. Este proceso proporcionó una comprensión detallada de la distribución de presiones a lo largo de la superficie del ala, así como los coeficientes de sustentación y resistencia. La optimización resultante permitió una eficiencia aerodinámica sobresaliente, fundamental para el rendimiento en vuelo. XFLR5 arrojo un desempeño para EHEKATL BWB entre 15 y 20 respecto a la polar $C_L/C_D$ en función de $\alpha$. En X-PLANE se logró un rendimiento entre 15 en las velocidades más bajas (Mach 0.04) y 27 en las velocidades más altas (Mach 0.085).
Una gran ventaja de nuestra metodología es que todos los software son de código abierto y gratuitos, a excepción de X-PLANE que para liberar el tiempo de simulación de quince minutos a ilimitado y todos los escenarios se tiene que comprar la versión completa, sin embargo, todo el trabajo aquí mostrado se desarrollo con la versión gratuita. 

El desarrollo del modelo dinámico de seis grados de libertad representó un hito crítico en la comprensión del comportamiento de la aeronave. Al considerar la dinámica longitudinal y lateral-direccional de manera independiente, se logró una representación precisa de los modos de vuelo individuales. Esto incluyó el desacoplamiento de las dinámicas no lineales, lo que facilitó un estudio detallado y específico de cada fase de vuelo.

La implementación de los controladores Backstepping y Adaptive Backstepping en ambos dominios, longitudinal y lateral-direccional, demostró ser altamente efectiva en la estabilización y control de la aeronave EHEKATL BWB. Estos controladores exhibieron una capacidad excepcional para adaptarse a condiciones cambiantes, lo que se tradujo en un control preciso y estable durante todas las fases de vuelo para su posterior implementación en los sistemas de control de vuelo.

Las simulaciones de vuelo en el entorno de X-PLANE brindaron una validación significativa de la efectividad de los controladores desarrollados. Las misiones de vuelo, que incluyeron despegue y ascenso, reconocimiento, así como descenso y aterrizaje, confirmaron la capacidad del EHEKATL BWB para operar con precisión y control en diversas condiciones y escenarios.



Este trabajo sienta las bases para una serie de investigaciones y desarrollos futuros, los cuales tienen el potencial de contribuir significativamente al campo de la aviación de vanguardia:

Explorar técnicas de optimización y ajuste de parámetros de controladores puede llevar a mejoras adicionales en la estabilidad y rendimiento en aeronaves de nueva generación, como es el caso de las aeronaves de ala-integrada, bajo una variedad de condiciones operativas y escenarios de vuelo. Esto incluye la adaptación de los controladores a situaciones específicas, como cambios en las condiciones atmosféricas.

La integración de sistemas avanzados, como aviónica de última generación y sistemas de propulsión innovadores, puede proporcionar un impulso adicional al rendimiento y capacidades operativas del EHEKATL BWB. La exploración de tecnologías emergentes en áreas como la propulsión eléctrica, de celdas de hidrógeno, o sistemas de control por asignación pueden abrir nuevas oportunidades en el diseño de aeronaves del futuro.

En perspectiva, este trabajo sienta las bases para futuras investigaciones y desarrollos en el campo de aeronaves de ala-integrada y sistemas de control de vuelo. La implementación exitosa de estrategias de control avanzadas como Backstepping y Backstepping Adaptable destaca el camino hacia aeronaves más eficientes, estables y seguras en el futuro de la aviación. La EHEKATL BWB emerge no solo como un logro técnico, sino como un símbolo del potencial mexicano en el panorama aeroespacial.



